\section{Related Work}
\label{sec:related}
\emph{Decoding motor imagery.} Common spatial patterns (CSP) find channel
combinations whose variance differs most between two
classes~\cite{ramoser2000,blankertz2008}. Riemannian methods classify each
trial's covariance matrix directly or after projecting it to a tangent
space~\cite{barachant2012,congedo2017}, and compact networks such as EEGNet
learn the filters from data~\cite{lawhern2018}. We use these standard
methods unchanged and ask how much patient data they need.

\emph{Shortening calibration.} Re-centring each person's covariance matrices
on their own mean makes data from different people
comparable~\cite{zanini2018}; Riemannian Procrustes analysis also matches
their spread and rotation~\cite{rodrigues2019}, and Euclidean alignment
applies re-centring to the raw trials, so that a network can use
it~\cite{he2020}. Such methods reduce calibration for healthy users of
conventional commands~\cite{lotte2015,wu2022review}. In stroke patients, transfer has mostly reused a
patient's own earlier sessions: it raised accuracy when few trials of the
current session were available~\cite{giles2022}, and transfer from earlier
sessions or from other patients helped the sessions that started with low
accuracy~\cite{cao2021}. We follow the whole calibration curve, up to the
full pool, for transfer from other stroke patients to a new command, with
re-centring only; Procrustes analysis is left untested.

\emph{Supernumerary limbs.} Extra robotic fingers and arms extend what a body
can do~\cite{prattichizzo2021,eden2022}, and wearing one changes the brain's
map of the hand~\cite{kieliba2021}. Sixth-finger imagery produced stronger
desynchronisation than imagery of a natural limb in healthy
people~\cite{wang2024jbhi}, and enhanced it in stroke
patients~\cite{wang2024tnsre}. We know of no study that asks how many trials
such a command needs, or whether it can be told apart from imagining the
hand.

\emph{Evaluation pitfalls.} Fitting any supervised step on data that later
serve as test data inflates accuracy~\cite{kapoor2023,varoquaux2017}. In
EEG, splitting segments of one person across training and test inflates
it~\cite{brookshire2024}, and block designs let decoders exploit slow
drifts~\cite{li2021perils}. We measure the size of one such leak on this
dataset and on two healthy datasets.
